% !TeX root = ../../thesis.tex

\subsection{\CI}
In traditional software development, the design phase usually leads to a representation of the required functionality in multiple, stand-alone modules. These modules are then implemented separately by the individual developers and are afterwards integrated into one monolithic application. This operation can prove to be very complicated since every developer can make their assumptions at the start of the project. Ultimately, these assumptions may render the components mutually incompatible. Furthermore, since it can take several weeks to months before this integration takes place, the developers often need to rewrite old code. Eventually, this will lead to unanticipated delays and costs \cite{SmartJenkinsDefinitive}.\\

\noindent Contrarily, agile development methodologies advocate the idea of frequent, yet small deliverables. In order to obtain frequent deliverables, we require to build the code often and integrate the modules multiple times, \emph{continuously}, rather than just once at the end. Another advantage of frequent deliverables is the early identification of problems \cite{gitlabci}. We refer to this practice as \acrfull{ci} \cite{SmartJenkinsDefinitive}. Note that this idea has existed before the creation of the Agile manifesto. The first notorious software company that has adopted CI is Microsoft in 1989 \cite{cusumanomicrosoft}. Cusumano reports that Microsoft typically builds the entire application at least once per day and as such requires developers to integrate and test their changes multiple times per day.\\

\noindent The introduction of \CI{} in software development has significant consequences on the life cycle. Where the waterfall model used a cascading life cycle, CI employs a circular, repetitive structure consisting of three phases, as visualised in \Cref{fig:agile-ci-lifecycle}.

\begin{enumerate}
	\bolditem{Implementation} In the first phase, every developer writes code individually for their assigned module. At a regular interval, the code is committed to the remote repository.
	
	\bolditem{Integration} When the developer commits their changes, they simultaneously fetch the changes other developers have made to their modules. Afterwards, the developer must integrate these remote changes with their local module to ensure these remain compatible. In case of a conflict, the developer is responsible for resolving this locally \cite{martin2014}.
	
	\bolditem{Test} After the developer has successfully integrated the remote modules to the locale module, the test suite must be executed to verify that no regressions have been introduced.
\end{enumerate}

\begin{figure}[htbp!]
	\centering
	\includegraphics[width=0.45\textwidth]{assets/images/ci-lifecycle.pdf}
	\caption{Development Life Cycle with \CI{}.}
	\label{fig:agile-ci-lifecycle}
\end{figure}

\noindent Adopting \CI{} can prove to be a lengthy and repetitive task. Luckily, a variety of tools and frameworks exist to automate this process. These tools are typically attached to a version control system (e.g. Git, Mercurial, \dots) using a \texttt{post-receive} hook \cite{SmartJenkinsDefinitive}. Every time the developers push a commit to the repository, the hook will notify the CI system. The CI system will respond accordingly by automatically building the code and executing the test cases. Optionally, we can configure the CI system to automatically publish successful builds to the end-users, a process referred to as \emph{Continuous Delivery}. This paper will now proceed by discussing four prominent CI systems.

% !TeX root = ../../../thesis.tex

\subsubsection{Jenkins}
Jenkins CI\footnote{\url{https://jenkins.io/}} was started as a hobby project in 2004 by Kohsuke Kawaguchi, a former employee of Sun Microsystems. Jenkins is programmed in Java and profiles itself as ``The leading open source automation server''. It was initially named Hudson, but after Sun was acquired by Oracle, issues related to the trademark Hudson arose. In response, the developer community decided to migrate the Hudson code to a new repository and rename the project to Jenkins \cite{SmartJenkinsDefinitive}. As of today, Jenkins is still widely used for many reasons. Since it is open source and its source code is located on GitHub, it is free to use and can be self-hosted by the developers in a private environment. Furthermore, Jenkins provides an open ecosystem to support developers into writing new plugins and extending its functionality. A market research conducted by ZeroTurnaround in 2016 revealed that Jenkins is the preferred Continuous Integration tool by 60\% of the developers \cite{maple_2016}.

\begin{figure}[htbp!]
	\centering
	\includegraphics[width=0.5\textwidth]{assets/jenkins.pdf}
	\caption{Logo of Jenkins CI (\url{https://jenkins.io/})}
	\label{fig:jenkins}
\end{figure}
% !TeX root = ../../../thesis.tex

\subsubsection{\github{} Actions}\label{sssec:github-actions}
Following the successful beta of \githubactions{} which had started in August 2019, \github launched its \CI{} system later that year in November\footnote{\url{https://github.blog/2019-08-08-github-actions-now-supports-ci-cd/}}. \githubactions{} executes builds in the cloud on servers owned by \github{}, restricting its use to \github{} repositories. Support for \github{} Enterprise repositories is not currently available. Developers can define builds using workflows, which can run both on Linux, Windows as well as macOS hosts. Private repositories are allowed a fixed amount of free build minutes per month, while builds of public repositories are always free of charges  \cite{githubactions}. Similar to Jenkins, we can extend \githubactions{} with custom plugins. These plugins can be composed either using a Docker container or using native JavaScript \cite{aboutgithubactions}. However, since this service is recent, it does not yet offer many plugins.

\begin{figure}[htbp!]
	\centering
	\includegraphics[width=0.5\textwidth]{assets/images/github-actions.pdf}
	\caption{Logo of \github{} Actions (\url{https://github.com/features/actions}).}
	\label{fig:github-actions}
\end{figure}
% !TeX root = ../../../thesis.tex

\subsubsection{GitLab CI}
GitLab, the main competitor of \github{}, announced its own \CI{} service in late 2012 named GitLab CI\footnote{\url{https://about.gitlab.com/blog/2012/11/13/continuous-integration-server-from-gitlab/}}. GitLab CI builds are configured using \emph{pipelines} and are executed by \emph{GitLab Runners}. These runners are operated by developers on their own infrastructure. Additionally, GitLab also offers the possibility to use \emph{shared runners}, which are hosted by themselves \cite{ciusinggitlab}. Equivalent to the aforementioned \github{} Actions, shared runners can be used for free by public repositories and are bounded by quota for private repositories \cite{gitlabdocs}. A downside of using GitLab CI is the lack of a community-driven plugin system, however this is a planned feature \footnote{\url{https://gitlab.com/gitlab-org/gitlab/issues/15067}}.

\begin{figure}[htbp!]
	\centering
	\includegraphics[width=0.40\textwidth]{assets/images/gitlab.pdf}
	\caption{Logo of GitLab CI (\url{https://gitlab.com/})}
	\label{fig:gitlab-ci}
\end{figure}
% !TeX root = ../../../thesis.tex

\subsubsection{\travisci}
\label{sssec:travisci}
The final CI platform that will be discussed is \travisci{}. This CI system was launched in 2011 and is only compatible with \github{} repositories. \travisci{} build tasks are configured similarly to GitLab CI but lack support for self-hosted runners. In addition to commit-triggered builds, we can also schedule daily, weekly or monthly builds using \emph{cronjobs}. Similar to \githubactions{}, open-source projects can use this service at zero cost, and a paid plan exists for private repositories \cite{travisanalysis}. It is not possible to create custom plugins, but \travisci{} already features built-in support for a variety of programming languages. In 2020, almost one million projects are using \travisci{} \cite{travis2020}.

\begin{figure}[htbp!]
	\centering
	\includegraphics[width=0.40\textwidth]{assets/images/travis-ci.pdf}
	\caption{Logo of \travisci{} (\url{https://travis-ci.com/}).}
	\label{fig:travis-ci}
\end{figure}
